\par\clearpage
\noindent{\small\textbf{ON THE HANOVER TRIANGULATION.}}\hfill359\par\medskip
\begin{center}
\textsc{Gauss} to \textsc{Bessel}.\quad Göttingen, 5~November 1823.
\end{center}
\begingroup
\makeatletter
\renewcommand{\@makefnmark}{\hbox{\normalfont[\textsuperscript{*})]}}
\long\def\@makefntext#1{\parindent1.5em\noindent[\textsuperscript{*})\ #1}
\makeatother

\par\noindent\hspace*{\parindent}.\ .\ .\ .\ .\ .\quad I have spent part of the year completing the remaining portion of my trigonometrical measurements in the north: from Timpenberg, Nindorf, and Lüneburg as far as Hamburg. Then I also made a preliminary reconnaissance of the area lying farther west, toward Bremen, since our government wishes the measurements to be extended farther west, as far as the Dutch border and the junction with \textsc{Krayenhoff}'s triangles. Finally I visited the Brocken and Hohehagen once more, partly because the angles of the triangle Hohehagen-Hils-Brocken had been measured under very unfavorable conditions in 1821, and partly because the connection of the Brocken with the Inselsberg, which had failed then, still had to be made, as did the connection to the Hessian triangulation point Meisner.
I have achieved these aims \textls[120]{for the most part} to my satisfaction. I would have liked to visit the Hils once more, too, in order to measure the angle there more accurately, but the lateness of the season prevented me.
Apart from this circumstance, I can now consider the triangulation for the arc measurement complete, insofar as it falls within my remit.
No further astronomical observations have been made beyond those needed for the orientation of my first triangle side.
It is still very uncertain, incidentally, whether \textls[120]{I} will carry out the above-mentioned plan to extend the measurements westward.
There is much in its favor and perhaps even more against it, even apart from the fact that a change in my external circumstances might yet become possible\footnote{This concerned the appointment to Berlin.]}.\par

\par\noindent\hspace*{\parindent}.\ .\ .\ .\ .\ .\quad In the past few days I have carefully adjusted the system of my principal triangles, so that the sum of the angles in each individual triangle and the ratios of the sides in the crossed quadrilaterals and pentagons agree exactly, all without arbitrariness, selection, or exclusion, and all with the rigor of the calculus of probabilities.
There are 26 triangles in all, in which I observed every angle myself.
\par\clearpage
\noindent360\hfill{\small\textbf{CORRESPONDENCE WITH BESSEL.}}\par\medskip
\noindent The greatest sum of the errors is \(2^{\prime\prime}{,}2\) in a triangle where sighting along one side was very difficult. The next greatest is \(1^{\prime\prime}{,}8\).
None of the 76 directions involved was changed by a full arcsecond in the adjustment. The greatest change amounts to \(0^{\prime\prime}{,}813\), on the side from Nindorf to Hamburg mentioned above.
What, according to my new theory of probability, I call the mean error is \(0^{\prime\prime}{,}48\) for the directions.
\par\endgroup

